\documentclass[10pt,a4paper]{amsart}
\usepackage[margin=19mm]{geometry}
\usepackage{amsmath,amssymb,mathtools}
\usepackage[scr=rsfs,cal=euler]{mathalfa}
\usepackage{xfrac}
\usepackage[unicode,hidelinks,bookmarks=false]{hyperref}
\newcommand{\ie}{i.e.\ }
\newcommand{\cf}{cf.\ }
\newcommand{\resp}{respectively\ }
\newcommand{\Subsection}{\subsection}
\providecommand{\nobreakdash}{\nobreak\hbox{-}}
\setlength{\emergencystretch}{2em}
\multlinegap0pt
\usepackage{tikz}
\usepackage{amssymb}
\usepackage[shortlabels]{enumitem}
\usepackage{enumerate}
\usepackage{mathtools}
\usepackage[normalem]{ulem}
\newtheorem{thm}{Theorem}
\newcommand*{\abs}[1]{\lvert #1\rvert}
\DeclareMathOperator{\irr}{irr}
\title{Maximum ratio of (graph) irregularities}
\author{Stijn Cambie, Jionghua Chang}
\date{Source: arXiv:2604.25341v1 (CC BY 4.0)}
\begin{document}
\maketitle

\noindent\emph{Source and licence.} Maximum ratio of (graph) irregularities. Version arXiv:2604.25341v1 (CC BY 4.0), distributed by arXiv under the Creative Commons Attribution 4.0 International licence, http://creativecommons.org/licenses/by/4.0/, as recorded in the arXiv licence metadata for this exact version and shown on the abstract page https://arxiv.org/abs/2604.25341v1. Full licence notice: \texttt{licenses/arxiv-2604.25341-CC-BY-4.0.txt}.\par\smallskip

\noindent\textit{Editorial scope.} Counted unit: the article's thm thm:conj25\_elegant (source0.tex lines 157--163). The article's complete original proof of it is delivered with it (source0.tex lines 165--195, one \texttt{proof} environment of 31 lines). No mathematics is written, repaired or re-derived: the statement and the proof are the article's own lines, in the article's own notation, and no retained line is substituted or re-flowed.  No further source-local context is needed: every cross-reference of every retained line resolves inside the two delivered spans.  Nothing was omitted from any retained span, so the package carries no omission record.  The retained lines cite 1 works, whose entries are transcribed verbatim from the article's own bibliography source in the e-print and delivered with short editorial labels recorded in the item's reference mapping.  No retained line contains a figure environment, an image-inclusion command or drawing code, so no image is delivered and the package declares no asset dependency.  Editorial formatting only, in addition: an AMS \texttt{amsart} preamble that restates the article's own declarations and macros used by the retained lines, namely the environments \texttt{thm} on separate counters, together with the standard packages those lines need.  The retained environments print in this document as Theorem 1 in order of appearance, whereas the article prints the same items with the numbering of its own full document.  The complete native source of the article as distributed in the arXiv e-print for this exact version is bundled unchanged as \texttt{sources/199296/source0.tex}.
\begin{thm}\label{thm:conj25_elegant}
    For every tree $T=(V, E)$ with order at least $3$,
    % $$n Var(T) <
    $$\sum_{e=uv \in E} \abs{d(u)-d(v)} \ge \sum_{v \in V} ( d(v)-2)^2 +2(\Delta(T)-2)  .$$
    % \le  \sum_{e=uv \in E} (d(u)-d(v))^2
    % Equality is only attained by paths.
\end{thm}

\begin{proof}
    Let $T$ be a tree which is a counterexample.
    By~\cite[Thm.~2.1]{DGLC23} (the proof works analogously for minimising $\irr(T)$ instead of $\sigma(T)$), we can assume it is the greedy tree with a certain degree sequence (the right hand side is fixed by the degree sequence, while the left hand side is not).
    Greedy trees are constructed from a given degree sequence by assigning the highest degree $\Delta$ to the root, the second-, third-, $(\Delta+1)^{th}$ highest degrees to the root's neighbours, and so on.



    We consider this greedy tree as a rooted tree with the root $r$ having maximum degree, drawn with the root on top and children of every vertex drawn below.
    % Now we consider the tree from leaves to root (or bottom to top).

    For every vertex $x$, let $T_x$ be the tree rooted in $x$ with its descendants (the vertex $x$ with everything connected to it below).
    
    % Still writing $d$ to denote $d_G$,
    We prove by induction on the height of $T_x$ that for every $x \not =r$,
    $$ g(T_x)=\sum_{e=uv \in E(T_x) } \abs{d(u)-d(v)} - \sum_{v \in V(T_x)} ( d(v)-2)^2 = d(x)-2.$$

    The base case is true, since for a leaf $x$ it states $0-1=1-2.$
    % For every leaf $v$, the contribution $(d(v)-2)^2$ equals $1.$ 
    
    If it is true for the children of vertex $x$,
    with children $Y$, then
    
\begin{align*}
    % \sum_{e=uv \in E(T_x) } \abs{d(u)-d(v)} - \sum_{v \in V(T_x)} ( d(v)-2)^2 & = \sum d(x)-d(y) -(d(x)-2)^2 +    \sum_{e=uv \in E(T_y) } \abs{d(u)-d(v)}- \sum_{v \in V(T_y)} ( d(v)-2)^2
    g(T_x)&= - (d(x)-2)^2+ \sum_{y \in Y} (d(x)-d(y))  + \sum_{y \in Y} g(T_y)\\&= - (d(x)-2)^2 +  \sum_{y \in Y} (d(x)-2) \\&= d(x)-2. 
\end{align*}

Finally, for the root $r$, the same computation leads to 
$$\hspace{2cm} g(T_r)= -(d(r)-2)^2 + d(r)(d(r)-2)= 2(d(r)-2)=2(\Delta(T)-2). \hspace{2cm} \qedhere$$
% The difference between the two quantities is thus at least $2(\Delta(T)-2),$ so equality is only valid if $\Delta(T)=2$ (noticing that it has order at least $3$ and thus $\Delta \ge 2$), i.e., $T$ is a path.
\end{proof}

\begin{thebibliography}{99}
\bibitem[DGLC23]{DGLC23} D.~Dimitrov, W.~Gao, W.~Lin, and J.~Chen. \newblock Extremal trees with fixed degree sequence for {{\(\sigma\)}}-irregularity. \newblock {\em DML, Discrete Math. Lett.}, 12:166--172, 2023.
\end{thebibliography}
\end{document}
